\documentclass[10pt,a4paper]{amsart}
\usepackage[margin=19mm]{geometry}
\usepackage{amsmath,amssymb,mathtools}
\usepackage[scr=rsfs,cal=euler]{mathalfa}
\usepackage{xfrac}
\usepackage[unicode,hidelinks,bookmarks=false]{hyperref}
\newcommand{\ie}{i.e.\ }
\newcommand{\cf}{cf.\ }
\newcommand{\resp}{respectively\ }
\newcommand{\Subsection}{\subsection}
\providecommand{\nobreakdash}{\nobreak\hbox{-}}
\setlength{\emergencystretch}{2em}
\multlinegap0pt
\usepackage{amssymb}
\usepackage{dsfont}
\usepackage{xcolor}
\usepackage[shortlabels]{enumitem}
\newtheorem{proposition}{Proposition}
\newtheorem{lemma}{Lemma}
\newcommand{\BA}{{\mathbb {A}}}
\newcommand{\BR}{{\mathbb {R}}}
\newcommand{\BZ}{{\mathbb {Z}}}
\newcommand{\CA}{{\mathcal {A}}}
\newcommand{\CM}{{\mathcal {M}}}
\newcommand{\XX}{\underline{\mathbf{X}}}
\title{Equidistribution of rational points and the geometric sieve for toric varieties}
\author{Zhizhong Huang}
\date{Source: arXiv:2111.01509v3 (CC BY 4.0)}
\begin{document}
\maketitle

\noindent\emph{Source and licence.} Equidistribution of rational points and the geometric sieve for toric varieties. Version arXiv:2111.01509v3 (CC BY 4.0), distributed by arXiv under the Creative Commons Attribution 4.0 International licence, http://creativecommons.org/licenses/by/4.0/, as recorded in the arXiv licence metadata for this exact version and shown on the abstract page https://arxiv.org/abs/2111.01509v3. Full licence notice: \texttt{licenses/arxiv-2111.01509-CC-BY-4.0.txt}.\par\smallskip

\noindent\textit{Editorial scope.} Counted unit: the article's proposition prop:subvar (source0.tex lines 1777--1784). The article's complete original proof of it is delivered with it (source0.tex lines 1815--1835, one \texttt{proof} environment of 21 lines, which the article places away from the statement). No mathematics is written, repaired or re-derived: the statement and the proof are the article's own lines, in the article's own notation, and no retained line is substituted or re-flowed.  Retained uncounted as the source-local context the proof needs: lines 1272--1274; lines 1802--1804; lines 1805--1807.  Nothing was omitted from any retained span, so the package carries no omission record.  No retained line contains a citation, so the wrapper carries no bibliography and no bibliography entry.  No retained line contains a figure environment, an image-inclusion command or drawing code, so no image is delivered and the package declares no asset dependency.  Editorial formatting only, in addition: an AMS \texttt{amsart} preamble that restates the article's own declarations and macros used by the retained lines, namely the environments \texttt{proposition}, \texttt{lemma} on separate counters, together with the standard packages those lines need.  The retained environments print in this document as Proposition 1, Lemma 1 in order of appearance, whereas the article prints the same items with the numbering of its own full document.  The complete native source of the article as distributed in the arXiv e-print for this exact version is bundled unchanged as \texttt{sources/117979/source0.tex}.
\begin{equation}\label{eq:CAsigmaB}
	\CA_\sigma(B):=\CA(B)\cap \pi^{-1}(C_{\sigma}(\BR)).
\end{equation} 

Before entering into the proof of Proposition \ref{prop:subvar}, we consider, for every $\sigma\in\triangle_{\max}$ and $1\leqslant k_0\leqslant n$, the set \begin{equation}\label{eq:CABik}
	\CA_\sigma(B)_{k_0,\widetilde{\eta}[B]}:=\{\XX\in\CA_\sigma(B):X_{k_0}=\widetilde{\eta}[B](X_k,1\leqslant k\leqslant n,k\neq k_0)\},
\end{equation} where $\widetilde{\eta}[B]:\BZ_{>0}^{n-1}\to\BR_{\geqslant 1}$ is a map depending on $B$.

\begin{lemma}\label{le:CABik}
		If $1\leqslant k_0\leqslant r$, assume moreover that $\widetilde{\eta}[B]$ depends only on $X_i,1\leqslant i\leqslant r,i\neq k_0$. We then have $$\#\CA_\sigma(B)_{k_0,\widetilde{\eta}[B]}=O(B(\log B)^{r-2}\log\log B).$$
\end{lemma}

\begin{proposition}\label{prop:subvar}
	Let $\phi$ be a non-zero integral polynomial in $n$-variables. Let us consider the set 
\begin{equation}\label{eq:CAphi}
		\CA_{\phi}(B):=\{\XX\in\CA(B):\phi(\XX)=0\}.
\end{equation}
	Then we have
	$$\#\CA_{\phi}(B)\ll (\deg \phi) B(\log B)^{r-2}\log\log B,$$ where the implied constant does not depend on $\phi$.
\end{proposition}

\begin{proof}[Proof of Proposition \ref{prop:subvar}]
	We recall the set  \eqref{eq:CAsigmaB}. For every $\sigma\in\triangle_{\max}$, we define $$\CA_{\sigma,\phi}(B):=\{\XX\in\CA_\sigma(B): \phi(\XX)=0\}.$$ It suffices to prove,  uniformly for $\phi$, 
\begin{equation}\label{eq:CAsigmaphi}
	\#\CA_{\sigma,\phi}(B)\ll (\deg \phi)B(\log B)^{r-2}\log\log B.
\end{equation}

	We fix an admissible ordering with respect to $\sigma$ and consider the projection $p_n:\BA^n\to\BA^{n-1}$ onto the first $n-1$ coordinates. Then either $\phi$ is constant on the variable $x_{n}$, i.e., $\phi$ factors through $p_n$, or the restriction $\pi_n:=p_n|_{\phi=0}$ is a generically finite dominant morphism to $\BA^{n-1}$. For the latter case, consider the locus $Y_{n-1}\subset\BA^{n-1}$ where $\pi_n$ fails to be finite. If we regard $\phi$ as a (non-constant) polynomial in $x_n$, then for any $\XX\in\BZ^n$ with $\phi(\XX)=0$, $\pi_n(\XX)\in Y_{n-1}$ if and only if all the coefficients (as polynomials in $\BZ[x_1,\cdots,x_{n-1}]$) of $\phi$ vanish on $\pi_n(\XX)$. We can take for example the (non-zero) polynomial $\phi_{n-1}\in\BZ[x_1,\cdots,x_{n-1}]$ among these coefficients of the highest degree, and we define
	$$\CA_{\sigma,\phi}^{(1)}(B):=\{\XX\in\CA_{\sigma,\phi}(B):p_n(\XX)\not\in Y_{n-1}\},$$
	$$\CA_{\sigma,\phi_{n-1}}(B):=\{\XX\in\CA_\sigma(B):\phi_{n-1}(p_n(\XX))=0\},$$
	so that $\deg\phi_{n-1}\leqslant \deg\phi$ and $$\CA_{\sigma,\phi}(B)\subset\CA_{\sigma,\phi}^{(1)}(B)\bigcup \CA_{\sigma,\phi_{n-1}}(B).$$ Note however for the first case where $\phi$ is constant in $x_n$ we just have  $\CA_{\sigma,\phi}(B)=\CA_{\sigma,\phi_{n-1}}(B)$ and $\CA_{\sigma,\phi}^{(1)}(B)=\varnothing$.
	Continuing the analysis for $\BA^{n-1}$ and the non-zero polynomial $\phi_{n-1}$ (if non-constant), we can further break $\CA_{\sigma,\phi_{n-1}}(B)$ into subsets. Finally we obtain 
	\begin{equation}\label{eq:decompCABf}
		\CA_{\sigma,\phi}(B)\subset \bigcup_{k=1}^{n}\CA_{\sigma,\phi}^{(k)}(B),
	\end{equation}
	where each set $\CA_{\sigma,\phi}^{(k)}(B),1\leqslant k\leqslant n$, if non-empty, is constructed  at the $k$-th step with respect to a non-constant polynomial $\phi_{n-k+1}\in\BZ[X_1,\cdots,X_{n-k+1}]$ with $\deg \phi_{n-k+1}\leqslant \deg \phi$ (by convention $\phi_n:=\phi$), and has the property that,  for $1\leqslant k\leqslant n-1$, there exists a finite collection of maps $(\eta_{k}(m):\BZ_{>0}^{n-k}\to\BR_{\geqslant 1})_{m\in \CM(k)}$, such that for every $\XX\in\CA_{\sigma,\phi}^{(k)}(B)$,  $X_{n-k+1}=\eta_{k}(m)(X_1,\cdots,X_{n-k})$ for a certain $m\in \CM(k)$, and $(\eta_{n}(m))_{m\in\CM(n)}$ consists of constant functions corresponding to the solutions of $\phi_1=0$ (in one variable). Moreover for every $1\leqslant k\leqslant n$,  $\#\CM(k)\leqslant\deg \phi_{n-k+1}\leqslant\deg \phi$.
	
	Recalling \eqref{eq:CABik}, by construction, the maps $(\eta_{k}(m):\BZ_{>0}^{n-k}\to\BR_{\geqslant 1})_{m\in \CM(k),1\leqslant k\leqslant n}$ meet the assumption of Lemma \ref{le:CABik}, and so we have $$\#\CA_{\sigma,\phi}^{(k)}(B)\leqslant \sum_{m\in\CM(k)} \#\CA_\sigma(B)_{n-k+1,\eta_k(m)}=O\left((\deg \phi)B(\log B)^{r-2}\log\log B\right).$$
	To conclude, we go back to \eqref{eq:decompCABf} and obtain
	$$\#\CA_{\sigma,\phi}(B)\leqslant \sum_{k=1}^{n}\#\CA_{\sigma,\phi}^{(k)}(B)=O\left((\deg \phi)B(\log B)^{r-2}\log\log B\right).$$
	The proof of \eqref{eq:CAsigmaphi} is thus completed.
\end{proof}

\end{document}
